% !TeX program = xelatex
\documentclass{gepa}

% Sur Overleaf : Menu > Compiler > XeLaTeX.
% Modifiez seulement les valeurs entre accolades.
% Si plusieurs personnes signent le travail, placez-les en ordre alphabétique de nom de famille.
\GEPATitre{Titre du travail}
\GEPASousTitre{Sous-titre facultatif}
\GEPAAuteur{Prénom nom}
\GEPADestinataire{Destiné à la professeur.e/chargé.e de cour Prénom Nom}
\GEPACours{Dans le cadre du cours POL000 Titre du cours}
\GEPALieu{Sherbrooke}
\GEPADate{Le date de remise}



% APA 7e édition en français sur Overleaf (Biber).
% La solution de repli sert seulement aux installations TeX locales minimales.
\IfFileExists{apa.bbx}{
  \usepackage[backend=biber,style=apa,sortlocale=fr_CA]{biblatex}
  \DeclareLanguageMapping{french}{french-apa}
  \addbibresource{references.bib}
}{
  \usepackage[round,authoryear]{natbib}
  \let\parencite\citep
  \let\textcite\citet
  \newcommand{\printbibliography}{\bibliographystyle{apalike}\bibliography{references}}
}

\begin{document}

\maketitlegepa

% Décommentez seulement les pages liminaires nécessaires, dans cet ordre.
% \geparesume{Insérez le résumé ici.}
% \geparemerciements{Insérez les remerciements ici.}
% Si une IA générative a contribué au contenu ou à la méthodologie du travail,
% remplissez « Declaration IA.tex », puis décommentez la ligne suivante.
% \gepadeclarationia{
J'ai utilisé \textbf{[outil, fournisseur et version]} pour \textbf{[nature de
l'assistance]} dans \textbf{[sections concernées]}. Je n'ai transmis aucune
donnée confidentielle, aucun renseignement personnel ni aucune donnée non
publiée. J'ai vérifié les résultats par \textbf{[méthodes de validation]}.

% Ajoutez ici les mesures de protection particulières, s'il y a lieu.
}

\gepatableofcontents
% \gepalistoftables
% \gepalistoffigures
% \gepaacronymes{\begin{description}\item[ÉPA] École de politique appliquée\end{description}}

\input{écris_ici}

% Selon le GÉPA, placez les annexes AVANT les références bibliographiques.
% \clearpage
% \appendix
% \section{Titre de l'annexe}
% Contenu de l'annexe.

\clearpage
\begingroup
\singlespacing
\RaggedRight
\printbibliography
\endgroup

\end{document}
